\documentclass[12pt]{article}
\usepackage{notes-project}
\usepackage{graphicx}
\usepackage{subcaption}
\usepackage{listings}

\graphicspath{{lec02/}}
\numberwithin{figure}{section}

\lstset{
  language=Python,
  basicstyle=\ttfamily\footnotesize,
  frame=single,
  framesep=6pt,
  xleftmargin=6pt,
  xrightmargin=6pt,
  breaklines=true,
  breakatwhitespace=true,
  showstringspaces=false,
  columns=fullflexible,
  keepspaces=true,
  keywordstyle=\color{blue!60!black},
  commentstyle=\color{gray},
  stringstyle=\color{green!40!black},
}

\begin{document}

\begin{center}
  {\LARGE\bfseries Lecture 2: Analytical Methods}\\[0.25em]
  {\large Introduction to Data Analysis Techniques}\\[0.15em]
\end{center}

\section{Lagrange Interpolation} 

\begin{dfn}[Interpolation]
Methods to find a smooth function $f(x)$ that passes through a given set of points $(x_0, y_0), (x_1, y_1), \ldots, (x_n, y_n)$ 
\end{dfn} 

\begin{thm}[Unique Existence of Lagrange Interpolation]{lagrange-existence-uniqueness}
For any set of $n+1$ distinct points $(x_0, y_0), (x_1, y_1), \ldots, (x_n, y_n)$, there exists a unique polynomial of degree at most $n$ that passes through all these points.
\end{thm} 

\begin{conv}[Lagrange Interpolation Method]
    Given a set of $n+1$ distinct points $(x_0, y_0), (x_1, y_1), \ldots, (x_n, y_n)$. We may construct the following system of equations: 
    \begin{align*}
        \left\{
        \begin{array}{cccc}
            a_0 + a_1 x_0 + a_2 x_0^2 + \ldots + a_n x_0^n &=& y_0 \\
            a_0 + a_1 x_1 + a_2 x_1^2 + \ldots + a_n x_1^n &=& y_1 \\
            \vdots & & \vdots \\
            a_0 + a_1 x_n + a_2 x_n^2 + \ldots + a_n x_n^n &=& y_n
        \end{array} \implies 
        \begin{bmatrix}
            1 & x_0 & x_0^2 & \ldots & x_0^n \\
            1 & x_1 & x_1^2 & \ldots & x_1^n \\
            \vdots & \vdots & \vdots & \ddots & \vdots \\
            1 & x_n & x_n^2 & \ldots & x_n^n
        \end{bmatrix}
        \begin{bmatrix}
            a_0 \\
            a_1 \\
            \vdots \\
            a_n
        \end{bmatrix}
        =
        \begin{bmatrix}
            y_0 \\
            y_1 \\
            \vdots \\
            y_n
        \end{bmatrix} 
        \right. 
    \end{align*}

    According to the \color{purple}{Theorem in Linear Algebra} \color{black}, which states that a system of $n$ linear equations with $n$ unknowns has a unique solution if $\det{A} \neq 0$, where $A$ is the coefficient matrix. 

    Since we ask all points in the set to be distinct, we can know that $\det{A} \neq 0$. Therefore, there exists a unique solution $\vec a$ 

    So, we can find a unique polynomial $f(x) = a_0 + a_1 x + a_2 x^2 + \ldots + a_n x^n$ that passes through all the given points. 
\end{conv}


\end{document}
